\documentclass[12pt]{article}

\usepackage{fancyhdr}
\usepackage{enumitem}
\usepackage{amsmath}
\usepackage{amsfonts}
\usepackage{amssymb}

\renewcommand{\vec}[1]{\mathbf{#1}}

\title{Homework 3}
\author{Satvik Prasad}
\date{}

\pagestyle{fancy}
\lhead{Satvik Prasad}
\rhead{MATH3012}

\begin{document}
\begin{enumerate}
  \item
    \begin{enumerate}
      \item For an arbitrary rook polynomial coefficient $r_k,$ we
        choose $k$ rows and $k$ columns to be occupied by a rook, and
        then for each row we choose a corresponding column for the
        rook to be placed at, in $k!$ ways. Hence, this yields
        \begin{align*}
          r_k = {3 \choose k}^2 k!,
        \end{align*}
        meaning the rook polynomial for the $3 \times 3$ chessboard is
        \begin{align*}
          R(x) = 1 + 9x + 18x^2 + 6x^3.
        \end{align*}

      \item Once again, the arbitrary rook number will be given by
        \begin{align*}
          r_k = {8 \choose k}^2 k!,
        \end{align*}
        and therefore the rook polynomial for an $8\times8$ chessboard will be
        \begin{align*}
          R(x) = \sum_{k=0}^{8}{8 \choose k}^2 k!\cdot x^k.
        \end{align*}
    \end{enumerate}

  \item
    \begin{enumerate}
      \item The $5$ candies we need to distribute are m\&ms,
        butterfingers, peanut butter cups, skittles, and snickers.
        Thus, let us construct a chessboard as follows:

        \begin{center}
          \begin{tabular}{c|c|c|c|c|c|}
            & A              & B              & C              & D
            & E              \\
            \hline
            M  &                & $\blacksquare$ &                &
            $\blacksquare$ &                \\
            \hline
            B  & $\blacksquare$ &                &                &
            &                \\
            \hline
            P  & $\blacksquare$ &                & $\blacksquare$ &
            & $\blacksquare$ \\
            \hline
            Sk &                & $\blacksquare$ &                &
            &                \\
            \hline
            Sn &                &                &                &
            & $\blacksquare$ \\
            \hline
          \end{tabular}
        \end{center}
        Placing 5 'rooks' on this board will mean if any individual
        chooses a specific candy, they can only get one candy, and no
        one else can get the same candy.

        \newpage

      \item By swapping the 3rd and 1st rows, and then swapping the
        2nd and 3rd columns, we yield the following board:
        \begin{center}
          \begin{tabular}{|c|c|c|c|c|}
            \hline
            $\blacksquare$ & $\blacksquare$ &                &
            & $\blacksquare$ \\
            \hline
            $\blacksquare$ &                &                &
            &                \\
            \hline
            &                & $\blacksquare$ & $\blacksquare$ &
            \\
            \hline
            &                & $\blacksquare$ &                &
            \\
            \hline
            &                &                &                &
            $\blacksquare$ \\
            \hline
          \end{tabular}
        \end{center}

        Placing rooks on black squares (disallowed positions), we get
        $r_0 = 1,$ and $r_1 = 8.$

        For $r_k,$ $k \geq 2,$ we can either choose to include the
        top-right black square, or not. If we place a rook on it, the
        remainder of our board $B'$ has rook polynomial
        \begin{align*}
          (1+ x) \cdot (1 + 3x + x^2) = 1 + 4x + 4x^2 + x^3,
        \end{align*}
        whereas if we don't place a rook on it, the remainder of our
        new board $B''$ has rook polynomial:
        \begin{align*}
          (1 + 3x + x^2)^2(1 + x) &= (1 + 9x^2 + x^4 + 6x + 6x^3 +
          2x^2)(1 + x) \\
          &= (1 + 6x + 11x^2 + 6x^3 + x^4)(1 + x) \\
          &= 1 + 7x + 17x^2 + 17x^3 + 7x^4 + x^5
        \end{align*}
        Because $r_k = r_{k -1, B'} + r_{k, B''},$ or equivalently,
        $$R(\bar C, x) = xR(B', x) + R(B'', x),$$ we can compute
        $r_k$ from the coefficients of
        \begin{align*}
          &x(1 + 4x + 4x^2 + x^3) + 1 + 7x + 17x^2 + 17x^3 + 7x^4 + x^5 \\
          &= 1 + 8x + 21x^2 + 21x^3 + 8x^4 + x^5,
        \end{align*}
        yielding a new rook polynomial given by:
        \begin{align*}
          R(\bar C, x) = 1 + 8x + 21x^2 + 21x^3 + 8x^4 + x^5.
        \end{align*}

      \item Let $c_1 := \text{`A gets the wrong snack'}$, $c_2 :=
        \text{`B gets the wrong snack'},$ and so on. Then, the number
        of possible distributions can be counted as:
        \begin{align*}
          N(\bar c_1 \bar c_2 \ldots \bar c_5) = N - S_1 + S_2 - \ldots -S_5.
        \end{align*}
        Notice that $N$ is the total number of ways to distribute $5$
        unique snacks among $5$ people, yielding $N = 5!$ ways.
        Additionally, $S_k$ will be the sum of all the ways to
        distribute $k$ snacks incorrectly. This can be counted by
        using $r_k$ to fix at least $k$ of the snacks incorrect, and
        then arrange between the $5 -k$ remaining rows and columns in
        $(5 - k)!$ ways. Hence,
        \begin{align*}
          S_k = r_k \cdot (5 -k)!,
        \end{align*}
        yielding the expression:
        \begin{align*}
          N(\bar c_1 \bar c_2 \ldots \bar c_5)  &= 5! - 8\cdot 4! +
          21 \cdot 3! - 21 \cdot 2! + 8 \cdot 1! - 1 \cdot 0! \\
          &= 120 - 192 + 126 - 42 + 8 - 1 = 19
        \end{align*}
        possible ways to distribute the snacks.
    \end{enumerate}
\end{enumerate}
\end{document}
